%------------------------------------
% DIPLOMARBEIT: Hauptdatei Präsentation
%     
% Stefan Grünwald (2000 - 2007)
% Michael Fladischer (2011)
%
% Das Dokument lässt sich am besten mit "latexmk" bauen:
% http://www.phys.psu.edu/~collins/software/latexmk-jcc/
%
% Die genutzten LaTeX-Packages sind alle Bestandteil der TeXlive Distribution die mit Debian (Squeeze) ausgeliefert werden.
% Für eine bequeme Installation unter Debian genügt es, folgende Zeile als "root" auszuführen:
%
% aptitude install -R latexmk texlive-base texlive-bibtex-extra texlive-binaries texlive-common texlive-fonts-extra texlive-fonts-recommended \
%  texlive-formats-extra texlive-generic-extra texlive-generic-recommended texlive-humanities texlive-lang-german texlive-latex-base texlive-latex-extra \
%  texlive-latex-recommended texlive-latex3 texlive-math-extra texlive-pictures texlive-plain-extra texlive-pstricks texlive-publishers texlive-science \
%  latex-beamer
%
% Alle .tex Dateien sind für den vim Editor ausgelegt und enthalten Modeline-Anweisungen zur besseren Darstellung.
% 
% Patches zur Verbesserung von Fehlern, zur Optimierung oder für neue Features können gerne an <michael@fladi.at> geschickt werden.
%------------------------------------

% Dokumentenklasse%{{{
\documentclass{beamer}
%}}}

% Generische Packages%{{{
\usepackage[ngerman]{babel}
\usepackage[utf8]{inputenc}                 % Laden der UTF8-Eingabekodierung
                                            % ACHTUNG: Dateien müssen UTF8 codiert sein!
\usepackage[T1]{fontenc}                    % EC codierte Fonts
\usepackage{ae}                             % PDF encoding Info
\usepackage{palatino}                       % Brotschrift: Palationo
\usepackage{helvet}                         % Serifenlose Schrift: Helvetica
\usepackage[left]{eurosym}                  % EURO links
\usepackage{paralist}                       % Beeinflussung von Listen
\usepackage{tabularx}                       % Austreiben auf eine bestimmte Breite
\usepackage{longtable}                      % Mehrseitige Tabellen
\usepackage{ltxtable}                       % 'longtable' und 'tabularx'
\usepackage{array}                          % Zeilenumbruch in 'p'-Spalten
% http://www.ctan.org/tex-archive/macros/latex/contrib/booktabs/
% \toprule, \midrule und \bottomrule
\usepackage{booktabs}
% http://www.ctan.org/tex-archive/macros/latex/required/tools/
\usepackage{dcolumn}
% http://www.ctan.org/tex-archive/macros/latex/contrib/blindtext/
% Erzeugt Blindtext welcher beim Prüfen des Layouts behilflich sein kann
\usepackage{blindtext}
% http://ctan.org/tex-archive/macros/latex/contrib/siunitx/
\usepackage{siunitx}
% http://www.ctan.org/tex-archive/macros/latex/required/graphics/
\usepackage{color}
% http://www.ctan.org/tex-archive/macros/latex/contrib/xcolor/
\usepackage{xcolor}
% http://www.ctan.org/tex-archive/macros/latex/contrib/wasysym/
\usepackage{wasysym}
\usepackage{beamerthemesplit}
%}}}

% Lokale Einstellungen%{{{
% vim: set tw=160:

\def\campusType{DIPLOM-ARBEIT}
\def\campusTitle{To be defined}
\def\campusAuthor{Michael Fladischer}
\def\campusPersKnz{0810319004}
\def\campusCity{Graz}
\def\campusDate{\today}
\def\campusSubject{}
\def\campusKeywords{PGP,x509,Firefox,XML-DSig}

%}}}

% Weiche für PDF/DVI%{{{
% http://www.ctan.org/tex-archive/macros/latex/contrib/oberdiek/
% Bitmaps stinken ja in der Druckstufe sowieso, besser TIKZ benutzen :-)
\usepackage{ifpdf}
% für Encapsulated-Postscript und PDFTeX Grafiken
\ifpdf
    \graphicspath{{images/pdf/}}
    \pdfcompresslevel=9
\else
    \graphicspath{{images/dvi/}}
\fi
%}}}

% Glossaries%{{{
% http://www.ctan.org/tex-archive/macros/latex/contrib/glossaries/
\usepackage[acronym,style=long,nonumberlist]{glossaries} 
\loadglsentries{misc/glossary}
\makeglossaries
%}}}

% TODO-Notes %\left\{ {
% http://www.ctan.org/tex-archive/macros/latex/contrib/todonotes/
\usepackage{todonotes} 
\presetkeys{todonotes}{inline}{}
%}}}

% Benutzerdefinierte Befehle%{{{
\newcommand{\xmlelem}[1]{\index{#1}\texttt{<#1/>}}
\newcommand{\xmlattr}[2]{\index{#1!#2}\textit{#2}}
\newcommand{\yes}{$\circledast$}
\newcommand{\no}{$\circleddash$}
\newcommand{\optional}{$\circledcirc$}
%\newcommand{\yes}{\rlap{\Cross}\Square}
%\newcommand{\no}{\Square}
%\newcommand{\optional}{\rlap{\FilledSmallCircle}\Square}
%\newcommand{\yes}{\rlap{\FilledSmallCircle}\Circle}
%\newcommand{\no}{\Circle}
%\newcommand{\optional}{\rlap{\SmallCircle}\Circle}
%\newcommand{\todo}{\marginpar{\textbf{TODO}}}
\renewcommand*{\glspostdescription}{}

%}}}

% Campus02 CI%{{{
\setbeamercovered{transparent}
\mode<presentation>
{  
    \usetheme{Ilmenau}
    \useinnertheme{circles}
    \usefonttheme[onlymath]{serif}
    \setbeamercovered{transparent}
    \setbeamertemplate{blocks}[rounded][shadow=true]
    \definecolor{campus02orange}{RGB}{255,251,173}
    \definecolor{campus02grey}{RGB}{80,80,80}

    \setbeamercolor{structure}{fg=campus02grey}

    \setbeamercolor{palette primary}{fg=campus02grey,bg=campus02orange!70}
    \setbeamercolor{palette secondary}{fg=campus02grey,bg=campus02orange!80}
    \setbeamercolor{palette tertiary}{fg=campus02grey,bg=campus02orange!90}
    \setbeamercolor{palette quaternary}{fg=campus02grey,bg=campus02orange}

    \setbeamercolor{titlelike}{parent=palette quaternary}

    \setbeamercolor{block title}{fg=campus02grey,bg=campus02orange}
    \setbeamercolor{block title alerted}{use=alerted text,fg=campus02grey,bg=alerted text.fg!75!bg}
    \setbeamercolor{block title example}{use=example text,fg=campus02grey,bg=example text.fg!75!bg}

    \setbeamercolor{block body}{parent=normal text,use=block title,bg=block title.bg!25!bg}
    \setbeamercolor{block body alerted}{parent=normal text,use=block title alerted,bg=block title alerted.bg!25!bg}
    \setbeamercolor{block body example}{parent=normal text,use=block title example,bg=block title example.bg!25!bg}

    \setbeamercolor{sidebar}{bg=campus02orange!70}

    \setbeamercolor{palette sidebar primary}{fg=campus02grey}
    \setbeamercolor{palette sidebar secondary}{fg=campus02grey!75}
    \setbeamercolor{palette sidebar tertiary}{fg=campus02grey!75}
    \setbeamercolor{palette sidebar quaternary}{fg=campus02grey}

    \setbeamercolor*{separation line}{}
    \setbeamercolor*{fine separation line}{}
}
\logo{\includegraphics[width=0.2\textwidth]{campus02}}
%}}}

% Präsentations-Header%{{{
\title{Präsentation \campusType}
\subtitle{\campusTitle}
\author{\campusAuthor}
\institute{CAMPUS02}
\date{\campusDate}
%}}}

% Folien%{{{
\begin{document}

\frame{
    \titlepage
}

\frame {
    \frametitle{Inhalt}
    \tableofcontents
}

\section{Ausgangssituation}
\frame {
    \frametitle{Ausgangssituation}
    \begin{itemize}
        \item Kryptographische Sicherheit auf Transportebene $\Rightarrow$ \gls*{https}
        \item Verschiedene Autoren verfassen Inhalte auf einer Website
        \item Autoren verfassen Inhalte auf verschiedenen Websites
        \item Inhalte können nicht mit kryptographischen Merkmalen versehen werden
    \end{itemize}
}

\section{Forschungsfragen}
\frame {
    \frametitle{Forschungsfragen}
    \begin{itemize}
        \item Wie können digitale Signaturen in \gls*{xhtml} integriert werden, ohne dass dieser W3C-Standard dabei verletzt wird?
        \item Welche Möglichkeiten bietet \gls*{xml-dsig} zur wahlfreien Selektion von Bereichen innerhalb eines XHTML-Dokuments?
        \item Wie können die öffentlichen Schlüssel zu den Signaturen verteilt werden? Welche der bestehenden Technologien eignen sich dafür am besten?
    \end{itemize}
}

\section{Aufbau}
\frame {
    \frametitle{Aufbau}
    \begin{itemize}
        \item Einführung in \gls*{xml-dsig}
        \item Verbindungen zwischen \gls*{xml-dsig} und \gls*{xhtml}
        \item Signaturverfahren für \gls*{xml-dsig}
        \item Anforderungen an ein Browser-AddOn (\gls*{firefox})
    \end{itemize}
}

\section{Erkenntnisse}
\frame {
    \frametitle{Erkenntnisse}
    \begin{itemize}
        \item \gls*{xml-dsig} bietet große Flexibilität beim Signieren von \gls*{xml}-Dokumenten mit \gls*{x509} oder \gls*{openpgp}
        \item Mehrere Ansätze für Signatur-Bindung\\
            Hervorzuheben sind:
            \begin{itemize}
                \item Mikro-Formate mit fixen Signaturvorgaben
                \item \xmlelem{link}-Tag für vollständige Ausnutzung von \gls*{xml-dsig}
            \end{itemize}
        \item Signaturverfahren:
            \begin{itemize}
                \item \gls*{x509}: für Unternehmerischen oder staatlichen Einsatz geeignet
                \item \gls*{openpgp}: für schnelle Adaption durch Privatbenutzer geeignet, höherer Grad an Automation möglich
            \end{itemize}
    \end{itemize}
}

\section{Conclusio}
\frame {
    \frametitle{Conclusio}
    \begin{itemize}
        \item Weder \gls*{openpgp} noch \gls*{x509} decken alle Erfordernisse vollständig ab
        \item Einsatz einer Abstraktionsschicht um beide Verfahren zu ermöglichen
        \item Definition des Verfahrens in der Signatur selbst
    \end{itemize}
}

\section{Ausblick}
\frame {
    \frametitle{Ausblick}
    \begin{itemize}
        \item Entwicklung eines Prototypen für ein Browser-AddOn in \gls*{firefox}
        \item Zumindest Unterstützung für \gls*{openpgp} oder \gls*{x509}
    \end{itemize}
}
\end{document}
%}}}

% vim:filetype=tex foldmethod=marker autoindent expandtab shiftwidth=4 tw=160:
% Local variables:
% mode: latex
% End:              
